\begin{abstract}
The rapidly growing literature on adaptive model merging (such as AdaMerging and PolyMerge) claims to achieve state-of-the-art multi-task performance without training costs by adjusting merging coefficients at test time.
In this paper, we critically audit the foundations of this paradigm and expose a severe, previously unreported \emph{Task Suite Bias}: current evaluation protocols rely on a single, arbitrary combination of visual classification datasets that masks fundamental limitations.
By systematically partitioning these tasks into five distinct suites of varying domain distance and representational conflict, we reveal that online Test-Time Adaptation (TTA) methods suffer from transductive overfitting to local stream noise in high-conflict regimes, lagging behind offline validated counterparts in simulation and collapsing catastrophically below the static Uniform baseline in physical weight-space deployments.
Conversely, we demonstrate that a simple, regularized offline alternative---Offline Few-Shot Validation Tuning (OFS-Tune) using continuous low-degree polynomial trajectories (specifically linear $d=1$ and quadratic $d=2$ configurations)---completely bypasses test-time compute and filters transductive noise.
OFS-Tune consistently outperforms unconstrained online TTA (AdaMerging), and matches or outperforms the complex online PolyMerge ($d=2$) across all suites in simulation (achieving up to $68.62\%$ accuracy in the high-conflict Suite B and reducing performance variance across random seeds).
Crucially, our independent physical weight-space validation on CNN weights reveals that in real deep neural networks, OFS-Tune outperforms online PolyMerge by up to $3.70\%$ and online AdaMerging by up to $4.20\%$ because it avoids both transductive stream-level overfitting and the privileged task-routing assumptions required for unsupervised online adaptation.
Our findings challenge the prevailing adaptive model-merging consensus, calling for a shift toward more rigorous multi-suite evaluation practices.
\end{abstract}
